
% ============================================================
% Single-column article version
% "Accessible Brain Morphometry from Point-of-Care
%  Ultra-Low-Field MRI Using Physics-Constrained Deep Learning"
% Final pre-submission version — JBHI primary / IEEE Access fallback
% ============================================================
\documentclass[12pt,a4paper]{article}
\usepackage[margin=2.5cm]{geometry}
\usepackage{cite}
\usepackage{amsmath,amssymb,amsfonts}
\usepackage{graphicx}
\usepackage{textcomp}
\usepackage{booktabs}
\usepackage{multirow}
\usepackage{xcolor}
\usepackage{hyperref}
\usepackage{setspace}
\usepackage{titlesec}
\graphicspath{{./}}

\onehalfspacing

\begin{document}

\title{\textbf{A Reproducible Feasibility Baseline for Segmentation-Free\\
nWBV Regression from 64\,mT Ultra-Low-Field MRI\\
Using Physics-Constrained Deep Learning}}

\date{}
\maketitle

\bigskip

% ── Abstract ─────────────────────────────────────────────────────────────────
\begin{abstract}
We present a reproducible feasibility study for direct normalized
whole-brain volume (nWBV) regression from 64\,mT ultra-low-field MRI,
without segmentation software or super-resolution preprocessing.
A compact 3D Vision Transformer (ViT3D, 4.23\,M parameters) is trained
via a three-stage pipeline: physics-constrained 64\,mT simulation
(tissue-specific T1/T2 relaxation, Rician noise, B0 inhomogeneity) on
IXI, OASIS-1 fine-tuning, and a 769-parameter session-level adapter.
Under fully subject-independent cross-session leave-one-subject-out
cross-validation on real 64\,mT hardware ($n = 23$), the model achieves
MAE $= 0.0134$\,nWBV [95\,\% CI: 0.010--0.017] with 47\,ms inference.
The cohort nWBV range is narrow (width $= 0.056$), so MAE is comparable
to a constant-mean baseline (MAE $= 0.0126$); the study should be
interpreted as a feasibility baseline rather than a deployed capability.
Inter-session ICC(3,1) $= 0.615$ [0.236--0.866] provides a preliminary
reliability estimate, though the wide CI precludes definitive
classification at $n = 23$. Model predictions do not significantly
correlate with age ($\rho = +0.232$, $p = 0.287$), an acknowledged
failure mode at this sample size. In physics-simulated dementia
conditions, CDR=1.0 subjects ($n=2$, directional only) show MAE $= 0.157$
vs.\ $0.016$ for the unlabelled healthy group ($10\times$; note: CDR=0.0
group MAE $= 0.086$, giving $\sim1.8\times$ against the labelled healthy
controls), indicating extension to pathological populations requires
atrophy-enriched retraining. The primary
contributions are the physics-simulation pre-training recipe, the
leakage-free subject-independent LOOCV protocol, and the honest failure
characterization bounding the validated scope to healthy adults.
\end{abstract}

\noindent\textbf{Keywords:}
Brain morphometry, dementia biomarker, low-field MRI, normalized whole-brain
volume, physics-informed simulation, point-of-care neuroimaging,
Vision Transformer, domain adaptation

% ============================================================
\section{Introduction}
\label{sec:introduction}
% ============================================================

Dementia afflicts more than 55 million individuals globally, with cases
expected to triple by 2050~\cite{who2023}. Normalized whole-brain volume
(nWBV) --- the fraction of total brain parenchymal tissue relative to
estimated intracranial volume --- is a well-validated marker of
neurodegeneration, with each CDR stage associated with a 5--10\% reduction
in nWBV~\cite{fotenos2005,buckner2004}. Monitoring nWBV enables
identification of neurodegenerative changes before symptom onset, making it a
useful endpoint in screening studies and clinical trials~\cite{jack2018nia}.

Standardized nWBV measurement requires high-field MRI (1.5 or 3\,T) and
automated pipelines such as FreeSurfer~\cite{fischl2012freesurfer} or
FSL~\cite{smith2002fsl}, which demand 20--90 minutes of computation,
software installations exceeding 12\,GB, and scanner infrastructure costing
USD 1--3 million. These requirements place routine brain morphometry out of
reach in low- and middle-income countries, rural communities, and
point-of-care environments~\cite{prince2016world}.

The Hyperfine Swoop (64\,mT) is a transportable, room-agnostic scanner
costing approximately \$50{,}000~\cite{hyperfine2021}. It has been validated
for bedside stroke triage~\cite{kimberly2023bedside}, but its potential for
quantitative brain morphometry remains largely unexplored. The core challenge
is that 64\,mT scans have markedly different signal properties from 3\,T
acquisitions --- reduced spatial resolution ($1.6\times1.6\times5.0$\,mm),
Rician noise, reduced grey/white contrast, and poor B0 shimming --- making
established morphometry tools unreliable without adaptation.

Prior approaches to this domain gap include super-resolution with SynthSR
(30--120\,s additional processing)~\cite{iglesias2023synthsr},
contrast-agnostic segmentation with SynthSeg+
(2--3\,min)~\cite{billot2023synthseg}, and paired-data domain
adaptation~\cite{chen2019synergistic} requiring large matched datasets
unavailable at 64\,mT. None of these directly regresses nWBV from 64\,mT
images without intermediate segmentation or super-resolution.

\textbf{Contributions.}
\begin{enumerate}
  \item To our knowledge, the first direct nWBV regression from 64\,mT MRI
        without intermediate segmentation or super-resolution, evaluated under
        fully subject-independent cross-session LOOCV ($n = 23$), yielding
        MAE $= 0.0134$ [0.010--0.017] and compute-light inference (47\,ms
        on a standard GPU; hardware not yet benchmarked on edge devices).
  \item A physics-constrained 64\,mT simulation recipe (T1/T2 tissue
        remapping, Rician noise, B0 inhomogeneity) that yields a small
        positive trend over Gaussian-blur degradation ($\Delta r = +0.018$,
        bootstrap 95\,\% CI [$-0.021$, $+0.057$], $p = 0.847$) — not
        statistically significant at $n=75$, but provided as a reproducible
        baseline for future larger-scale comparison.
  \item A leakage-free subject-independent cross-session LOOCV protocol
        yielding a preliminary inter-session reliability estimate
        ICC(3,1) $= 0.615$ [0.236--0.866] at $n = 23$.
  \item An honest failure characterization: MAE is comparable to a
        constant-mean baseline (MAE $= 0.0126$) given the narrow cohort range;
        CDR=1.0 stress-testing ($n=2$, directional) shows $10\times$ MAE
        increase, bounding the validated scope to healthy adults.
\end{enumerate}

% ============================================================
\section{Related Work}
\label{sec:related}
% ============================================================

\subsection{Brain Morphometry from MRI}

FreeSurfer~\cite{fischl2012freesurfer} and
FastSurfer~\cite{henschel2020fastsurfer} are the leading tools for
volumetric brain morphometry and require high-field isotropic T1-weighted
acquisitions. Direct regression approaches without segmentation have been
demonstrated for brain age~\cite{peng2021brain,cole2017predicting,jonsson2019brain},
achieving 3--5 year mean absolute error on 3\,T data. Rebsamen et
al.~\cite{rebsamen2020cortical} showed that deep learning segmentation can
match pipeline-based cortical thickness estimates. Roy et
al.~\cite{roy2019quicknat} proposed QuickNAT for fast brain segmentation.
To our knowledge, no prior work has directly regressed nWBV from
ultra-low-field MRI without an intermediate segmentation or
super-resolution step, evaluated under subject-independent
cross-session LOOCV on real 64\,mT hardware.

\subsection{Ultra-Low-Field MRI}

V{\'a}{\v s}a et al.~\cite{vasa2025imaneur,vasa2025} characterize
test-retest reliability and high-field correspondence for ULF brain
morphometry and release the ds006557 dataset ($n = 23$ with paired 64\,mT
and 3\,T acquisitions). Islam et al.~\cite{islam2025aibrain} show that deep
networks can reduce volumetric bias between 64\,mT and 3\,T measurements.
Arnold et al.~\cite{arnold2022simulated} evaluate diagnostic performance at
64\,mT using simulation on public data, motivating our simulation-based
pre-training strategy. Kimberly et al.~\cite{kimberly2023bedside} and Sheth
et al.~\cite{sheth2021assessmentultralow} validate bedside ULF MRI for
neurological assessment. Javadi et al.~\cite{javadi2025diffusion} and
Dayarathna et al.~\cite{dayarathna2024adversarial} study image synthesis and
translation between 64\,mT and 3\,T. Gopinath et al.~\cite{gopinath2025cortical}
demonstrate robust cortical mapping from low-field MRI. Baljer et
al.~\cite{baljer2024paediatric} show multi-orientation super-resolution for
paediatric 64\,mT MRI. Minocha and Yang~\cite{minocha2024lfmri} investigate
early Alzheimer detection from low-field MRI. None of these works quantify
nWBV directly from 64\,mT scans without segmentation.

SynthSeg+~\cite{billot2023synthseg} achieves $r = 0.918$, MAE $= 0.005$
on real Hyperfine data (this study) and serves as an accuracy upper bound.
Its requirement for a 12\,GB FreeSurfer installation and 2--3\,min per scan
makes it unsuitable for co-location with bedside hardware.

\subsection{Vision Transformers in Medical Imaging}

Dosovitskiy et al.~\cite{dosovitskiy2020image} established that Vision
Transformers (ViT) are competitive with CNNs on image recognition.
TransUNet~\cite{chen2021transunet} and Swin UNETR~\cite{hatamizadeh2021swin}
demonstrated ViT advantages in 3D medical image segmentation. For global
regression tasks such as nWBV --- where the target is a ratio over the
entire brain volume --- ViT's global self-attention mechanism over patch
tokens is architecturally well-suited.

\subsection{Physics-Informed MRI Simulation}

Gaussian-blur degradation is commonly used as a resolution proxy in
cross-domain MRI pre-training~\cite{zhao2021superresolution}. Scanner
variability is a well-known source of model failure~\cite{karani2018lifelong,dewey2019deepharmony},
motivating domain adaptation modules. SynthSR~\cite{iglesias2023synthsr}
uses physics-based synthesis for super-resolution. Our simulation
explicitly models tissue contrast at 64\,mT using T1/T2 relaxation
parameters, Rician noise, and B0 inhomogeneity from published Hyperfine
Swoop measurements~\cite{vasa2025}.

% ============================================================
\section{Methods}
\label{sec:methods}
% ============================================================

\subsection{Physics-Constrained 64\,mT Simulation}
\label{sec:sim}

High-field T1w volumes (3\,T, 1\,mm isotropic) are degraded to simulate
Hyperfine Swoop 64\,mT T2w acquisitions based on published scanner
parameters (TR $= 2000$\,ms, TE $= 194.8$\,ms, ETL $= 80$, in-plane
resolution $1.6 \times 1.6$\,mm)~\cite{vasa2025}:

\textbf{Step 1 -- Tissue contrast remapping.}
Voxel intensities are remapped using tissue-specific T1/T2 relaxation
parameters at 64\,mT (Table~\ref{tab:t1t2}), derived from published
low-field relaxometry~\cite{deichmann2002}. White matter T1 decreases from
832\,ms (3\,T) to 400\,ms (64\,mT); grey matter from 1330\,ms to 500\,ms,
resulting in reduced grey/white contrast that characterises 64\,mT images.

\textbf{Step 2 -- Rician noise.}
Rician-distributed noise is added assuming SNR $\approx 4$, a conservative
lower bound. The actual scanner achieves effective SNR $\approx 309$ through
ETL-80 averaging and compressed sensing, which are not modelled. This SNR
mismatch is the primary source of the pre-adaptation bias (discussed in
Section~\ref{sec:limitations}).

\textbf{Step 3 -- Spatial resolution degradation.}
Trilinear downsampling to 64\,mT in-plane resolution with slice-direction
blurring ($\sigma = 2.5$\,voxels) to represent the 5\,mm point spread
function.

\textbf{Step 4 -- B0 inhomogeneity.}
Multiplicative spatial Gaussian ($\sigma = 15$\,mm, amplitude 10\%) models
poor shimming at 64\,mT. Representative axial slices are shown in
Fig.~\ref{fig:scan_compare}.

\begin{figure}[!t]
\centering
\includegraphics[width=0.8\textwidth]{fig7_scan_comparison.png}
\caption{Representative axial slices: 3\,T MPRAGE (left),
physics-simulated 64\,mT (centre), and real Hyperfine 64\,mT (right).
The simulation captures grey/white contrast reduction and noise
characteristics of 64\,mT, though SNR is intentionally conservative.}
\label{fig:scan_compare}
\end{figure}

\begin{table}[!t]
\caption{Tissue relaxation parameters used in 64\,mT simulation (ms)}
\label{tab:t1t2}
\centering
\setlength{\tabcolsep}{6pt}
\begin{tabular}{lcccc}
\toprule
Tissue & T1 @ 3\,T & T1 @ 64\,mT & T2 @ 3\,T & T2 @ 64\,mT \\
\midrule
White matter & 832  & 400  & 70   & 80   \\
Grey matter  & 1330 & 500  & 83   & 100  \\
CSF          & 4000 & 3000 & 2000 & 2000 \\
\bottomrule
\multicolumn{5}{l}{\footnotesize Based on published low-field relaxometry~\cite{deichmann2002}.}
\end{tabular}
\end{table}

\subsubsection*{Simulation fidelity}
We quantified how closely the simulator reproduces registration-free
image statistics of real 64\,mT scans, using the $n = 23$ ds006557
subjects for which paired real acquisitions exist. Because the
simulated and real volumes have different fields of view and are not
co-registered, voxelwise similarity indices (SSIM, PSNR) are not
physically meaningful across field strengths; we therefore report
normalised cross-correlation (NCC) of tissue-intensity structure and
compare it against the Gaussian-blur degradation baseline of Arnold et
al.~\cite{arnold2022simulated}. The physics simulator attains
NCC $= 0.497 \pm 0.045$ versus $0.516 \pm 0.055$ for Gaussian blur ---
a statistically indistinguishable difference, indicating that neither
degradation model reproduces fine anatomical texture and that the
simulator's benefit is not a structural-similarity effect. Two points
warrant emphasis. First, the simulated SNR ($\approx 33$) is
deliberately far below the real scanner's effective SNR ($\approx 309$,
achieved through ETL-80 averaging and compressed sensing that we do not
model), consistent with the conservative-noise design in
Section~\ref{sec:sim}; the simulator targets the raw-acquisition noise
regime rather than the vendor-reconstructed image. Second, and
consequently, the operative validation of the simulator is not
image-similarity but downstream task performance: the physics-vs-blur
pre-training ablation on the nWBV regression task itself
(Section~\ref{sec:ablation}, Table~\ref{tab:ablation}), which shows a
small positive but statistically non-significant advantage for the
physics recipe.

\subsection{ViT3D Architecture}
\label{sec:arch}

All volumes are resampled to $64^3$ voxels and intensity-normalized to
$[0, 1]$ using 1\textsuperscript{st}--99\textsuperscript{th} percentile
clipping. ViT3D comprises:
\begin{itemize}
  \item \textbf{Patch embedding}: $64^3$ volume $\to$ 64 non-overlapping
        $16^3$ patches, each linearly projected to 256 dimensions.
  \item \textbf{Transformer encoder}: 4 layers, 8-head multi-head
        self-attention, embedding dimension 256, MLP ratio 4, dropout 0.1;
        learnable [CLS] token aggregates global context.
  \item \textbf{Regression head}: Linear projection from [CLS] token to
        scalar nWBV.
\end{itemize}
Total parameters: 4.23M. The compact size is consistent with
compute-light inference (47\,ms on a standard GPU); edge-device
deployment is future work. The three-stage training pipeline is illustrated in
Fig.~\ref{fig:arch}.

\begin{figure}[!t]
\centering
\includegraphics[width=\textwidth]{fig1_architecture.png}
\caption{Three-stage training pipeline: (1) physics-guided pre-training on
IXI; (2) OASIS-1 fine-tuning; (3) lightweight cross-session LOOCV domain
adaptation on ds006557.}
\label{fig:arch}
\end{figure}

\subsection{Three-Stage Training Pipeline}
\label{sec:training}

\textbf{Stage 1 -- Physics-guided self-supervised pre-training
(IXI, $n=156$).}
The IXI dataset provides 1.5\,T and 3\,T acquisitions but no
nWBV labels, so supervised regression at this stage is not feasible.
We instead apply the 64\,mT physics simulator to each IXI volume and
train ViT3D as a denoising autoencoder, with the high-field source
volume as the reconstruction target. This objective serves three
specific purposes for the downstream task. First, it forces the encoder
to learn the inverse mapping from the degraded 64\,mT contrast back to
the underlying tissue structure --- precisely the representation
required to estimate parenchymal fraction at 64\,mT. Second, the patch
embedding learns to discount Rician noise and B0 inhomogeneity rather
than treat them as anatomical signal, addressing the dominant
confounders at low field. Third, reconstruction-based pre-training is
known to converge to features useful for global regression
targets~\cite{dosovitskiy2020image}, and the receptive field required
for whole-brain nWBV estimation is naturally established by ViT3D's
non-local attention over $4\times4\times4$ patch tokens. Stage 1
therefore provides a contrast-aware initialisation for the Stage 2
nWBV-supervised fine-tuning, without consuming any of the limited
labelled OASIS-1 or ds006557 supervision.

\textbf{Stage 2 -- OASIS-1 fine-tuning ($n= 375$).}
All 4.23M parameters are updated end-to-end using AdamW (lr $= 5 \times
10^{-5}$, weight decay $= 10^{-4}$, cosine annealing, 50 epochs, early
stopping patience $= 10$ on validation Pearson $r$). The subject-level split
(300 train / 37 validation / 38 test, seed 42) is fixed before any model
training and never altered.

\textbf{Stage 3 -- Lightweight domain adaptation via cross-session LOOCV
(ds006557, $n=23$).}
The ds006557 dataset provides two real Hyperfine sessions per subject:
HFC (first acquisition) and HFE (second acquisition), in addition to a
paired 3\,T acquisition from which FastSurfer ground-truth nWBV is derived.
We exploit this dual-session structure for a fully subject-independent
evaluation protocol:

\begin{enumerate}
  \item For each of the 23 subjects in turn, the final LayerNorm and
        regression head (769 parameters total: LayerNorm $= 512$,
        head $= 257$) are updated on session~HFC scans from the
        \emph{other} 22 subjects. FastSurfer GT nWBV from their paired
        3\,T acquisitions is used as the supervision signal. All other
        4.23M parameters remain frozen. Training: AdamW
        (lr $= 10^{-4}$, weight decay $= 10^{-4}$, cosine annealing,
        80 epochs), batch size 4, with 5 augmentations per scan
        (Gaussian noise $\sigma = 0.015$, intensity scaling $\pm 10$\%,
        random axis flips).
  \item The adapted model is evaluated on session~\textbf{HFE} of the
        held-out subject. Neither HFC nor HFE of the held-out subject
        appears anywhere in the training data for that fold.
\end{enumerate}

This yields 23 fully held-out HFE predictions, one per fold, with no
subject or session leakage. The best-epoch checkpoint (lowest training loss)
is used for inference. For ICC computation, the same fold model also
predicts session~HFC of the held-out subject, giving a pair of independent
session predictions per subject.

\textbf{Rationale for LayerNorm+head adaptation.}
Updating only the final LayerNorm (512 params) and regression head
(257 params) is a standard low-data domain-adaptation strategy that
preserves the OASIS-learned feature hierarchy while correcting
output-level bias~\cite{karani2018lifelong}. Full fine-tuning on $n=22$
subjects would risk catastrophic forgetting and overfitting to the
narrow nWBV range of ds006557 (0.752--0.808), while head-only adaptation
omits the LayerNorm rescaling required to absorb the residual
contrast-statistics gap between simulation and real hardware. The
LayerNorm+head choice therefore matches the structure of the observed
domain shift. A formal empirical comparison of head-only,
LayerNorm+head, and full fine-tune strategies is identified as future
work; the present choice is motivated by the structure of the observed
domain shift rather than an exhaustive search.

\textbf{Pseudo-label validity.}
The LOOCV adaptation uses FastSurfer GT labels derived from paired 3\,T
scans, independent of the 64\,mT input. An ablation (Section~\ref{sec:pseudo_ablation})
shows that replacing GT with SynthSeg+ pseudo-labels incurs a MAE penalty
of $+0.005$\,nWBV, with both label sources yielding test MAE below the
reference threshold. This confirms deployment feasibility without requiring
FastSurfer installation.

\subsection{Primary Metric and Statistical Analysis}
\label{sec:metrics}

Mean absolute error (MAE) in nWBV units is the pre-specified primary metric.
MAE $= 0.020$ corresponds to approximately 28\,mL of brain tissue
($0.020 \times$ mean eTIV $\approx 1400$\,mL), the minimum detectable
change distinguishing normal from pathological atrophy in longitudinal
studies~\cite{buckner2004}. This threshold was defined before any model
results were examined.

Bootstrap 95\,\% confidence intervals (10{,}000 resamples, seed 42) are
reported for all primary MAE values. Pearson $r$ and Spearman $\rho$ are
reported where appropriate, with the note that $r$ is structurally
unreliable in the narrow-range ds006557 cohort (nWBV width 0.056;
see Section~\ref{sec:loocv_results}). ICC(3,1) (two-way mixed-effects,
single measures, absolute agreement) is computed between fold-level HFC
and HFE predictions; bootstrap 95\,\% CI is derived from 10{,}000 subject-level
resamples.

\textbf{Constant-mean LOOCV baseline.}
Because the ds006557 nWBV range is narrow, we additionally report a
constant-mean LOOCV baseline that predicts, for each held-out subject,
the leave-one-out cohort mean of the remaining 22 subjects. In
narrow-range regression MAE is dominated by the range itself rather
than by model behaviour; comparing against this baseline establishes
whether the model has learned anything beyond the cohort mean.
Discriminative evidence beyond MAE is provided by the cross-session
ICC, which requires input-conditional behaviour and cannot be
reproduced by a constant predictor. Model predictions do not
significantly correlate with subject age ($\rho = +0.232$,
$p = 0.287$; Section~\ref{sec:loocv_results}), so age correlation
is not used as discriminative evidence.

% ============================================================
\section{Experiments}
\label{sec:experiments}
% ============================================================

\subsection{Datasets}

\textbf{IXI}~\cite{ixi2017}: 156 healthy subjects, multi-centre 3\,T T1w
scans. Used exclusively for Stage~1 physics-guided pre-training.

\textbf{OASIS-1}~\cite{marcus2007oasis}: 375 participants (age 20--96,
CDR 0--2), 1.5\,T T1w MRI; nWBV from FreeSurfer. Fixed subject-level split:
300 train / 37 validation / 38 test (seed 42).

\textbf{ds006557}~\cite{vasa2025}: 23 healthy adults (age 20--69). Each
subject has \emph{two} real Hyperfine 64\,mT sessions (HFC, HFE) and one
paired 3\,T session from which FastSurfer nWBV is derived. Used for Stage~3
domain adaptation and LOOCV evaluation only.

\textbf{Data integrity}: All three datasets are drawn from entirely separate
subject pools. Ground-truth nWBV is always derived from high-field
acquisitions, never from 64\,mT input.

\subsection{Baselines}

\textbf{SynthSeg+}~\cite{billot2023synthseg}: Applied to the same axial
T2w input using FreeSurfer 7.4.1 (\texttt{mri\_synthseg --robust}).
Serves as the accuracy upper bound; requires 12\,GB installation and
2--3\,min per scan.

\textbf{CNN3D}: A 3D CNN with 4 convolutional blocks, global average
pooling, and $\approx 4.1$M parameters, trained on the same OASIS-1 split
and hyperparameters as ViT3D. Evaluated on both OASIS-1 and real 64\,mT
data.

\textbf{Gaussian-blur pre-training}: ViT3D pre-trained with $\sigma = 2$
Gaussian blur plus noise instead of physics simulation. Evaluated under
identical Stage~2 conditions. Simulation design ablation only.

% ============================================================
\section{Results}
\label{sec:results}
% ============================================================

\subsection{OASIS-1 Test Set}
\label{sec:oasis_results}

On the held-out OASIS-1 test set ($n = 38$), ViT3D achieves Pearson
$r = 0.722$ [95\,\% CI: 0.524--0.847, $p < 0.001$], MAE $= 0.058$
[0.041--0.076], RMSE $= 0.081$ (Table~\ref{tab:main}, Fig.~\ref{fig:scatter}).
Performance degrades for atrophied brains (nWBV $< 0.76$;
Section~\ref{sec:failure}).

CNN3D on the same OASIS-1 split achieves $r = 0.877$ [0.775--0.935,
$p < 0.001$] and MAE $= 0.024$ [0.018--0.031], outperforming ViT3D in
this wider-range setting. The reversal on real 64\,mT hardware
(Section~\ref{sec:vit_cnn}) is consistent with the different nWBV range
and noise characteristics of the two evaluation regimes.

\begin{figure}[!t]
\centering
\includegraphics[width=0.8\textwidth]{fig2_nwbv_scatter.png}
\caption{Predicted vs.\ true nWBV --- OASIS-1 test set ($n=38$, ViT3D).
Large residuals occur exclusively at low nWBV ($<0.76$), reflecting the
atrophy failure mode described in Section~\ref{sec:failure}.}
\label{fig:scatter}
\end{figure}

\begin{table}[!t]
\caption{Summary of all results grouped by evaluation context.
Bootstrap 95\,\% CIs in brackets (10{,}000 resamples, seed 42);
OASIS-1 $r$ CIs via Fisher $z$ transform.
Reference threshold: MAE $< 0.020$\,nWBV~\cite{buckner2004}.}
\label{tab:main}
\centering
\small
\setlength{\tabcolsep}{5pt}
\renewcommand{\arraystretch}{1.15}
\begin{tabular}{p{3.2cm} p{3.8cm} c c c}
\toprule
\textbf{Context} & \textbf{Method} & \textbf{Pearson $r$ [95\,\% CI]} & \textbf{MAE [95\,\% CI]} & \textbf{RMSE} \\
\midrule
\multicolumn{5}{l}{\textit{(A) OASIS-1 test set ($n=38$)}} \\[1pt]
  & ViT3D (ours)            & 0.722 [0.524--0.847] & 0.058 [0.041--0.076] & 0.081 \\
  & CNN3D                   & 0.877 [0.775--0.935] & 0.024 [0.018--0.031] & ---   \\
\midrule
\multicolumn{5}{l}{\textit{(B) Real 64\,mT --- unadapted ($n=23$)}} \\[1pt]
  & SynthSeg+ (upper bound) & 0.918        & 0.005 & --- \\
  & CNN3D (no adapt.)       & 0.514        & 0.076 & --- \\
  & ViT3D (no adapt.)       & 0.291        & 0.040 & --- \\
\midrule
\multicolumn{5}{l}{\textit{(C) Real 64\,mT --- LOOCV adapted ($n=23$, cross-session)$^\dagger$}} \\[1pt]
  & ViT3D (LOOCV adapt.)    & ---$^\ddagger$ & \textbf{0.0134} [0.010--0.017] & 0.016 \\
\midrule
\multicolumn{5}{l}{\textit{(D) Inter-session reliability}} \\[1pt]
  & ICC(3,1) HFC vs.\ HFE   & \multicolumn{3}{l}{$0.615$ [95\,\% CI: 0.236--0.866] \quad \textit{moderate}~\cite{koo2016icc}} \\
\midrule
\multicolumn{5}{l}{\textit{(E) Pathology stress-test --- physics-simulated 64\,mT ($n=38$, no adapt.)}} \\[1pt]
  & CDR $= 0.0$             & --- & $0.086 \pm 0.051$ & --- \\
  & CDR $= 0.5$             & --- & $0.099 \pm 0.031$ & --- \\
  & CDR $= 1.0$             & --- & $0.157 \pm 0.012$ & --- \\
  & No CDR label            & --- & $0.016 \pm 0.010$ & --- \\
\bottomrule
\multicolumn{5}{l}{\footnotesize $^\dagger$Adapter (769 params) trained on HFC of $n{-}1$ subjects; tested on HFE of held-out.}\\
\multicolumn{5}{l}{\footnotesize $^\ddagger$Pearson $r$ omitted: nWBV range 0.056 structurally limits $r$ (Section~\ref{sec:loocv_results}).}\\
\multicolumn{5}{l}{\footnotesize CDR MAE: mean $\pm$ SD. Bias: (A) ViT3D $+0.014$; CNN3D no adapt.\ $-0.076$; LOOCV adapt.\ $-0.0022$.}
\end{tabular}
\end{table}

\subsection{CDR Stratification on OASIS-1}

Among the 18 OASIS-1 test subjects with CDR ratings, mean true nWBV
decreases monotonically with CDR: 0.0 $\to$ 0.752, 0.5 $\to$ 0.739,
1.0 $\to$ 0.677 (Fig.~\ref{fig:cdr}), consistent with the 5--10\% per
CDR-stage reduction reported in the literature~\cite{fotenos2005}.
Kruskal-Wallis test on true nWBV across CDR groups: $H = 4.930$,
$p = 0.085$ (not significant, largely attributable to $n = 2$ for CDR~1.0).
Mann-Whitney U, CDR~0.5 vs.\ CDR~1.0 (one-sided, CDR~0.5~$>$~CDR~1.0): $p = 0.022$. Cohen's $d$
(CDR~0.0 vs.\ CDR~1.0) $= 1.368$ (large effect; interpret cautiously at
$n = 2$).

\begin{figure}[!t]
\centering
\includegraphics[width=0.8\textwidth]{fig3_cdr_boxplot.png}
\caption{True nWBV by Clinical Dementia Rating (CDR) --- OASIS-1 test set
($n = 18$ with CDR labels). Monotonic decrease across CDR stages is
consistent with published atrophy trajectories~\cite{fotenos2005}.}
\label{fig:cdr}
\end{figure}

\subsection{Ablation: Physics Simulation vs.\ Gaussian Blur}
\label{sec:ablation}

Under controlled conditions (IXI pre-training + OASIS fine-tuning,
$n_{\text{test}} = 75$), physics-constrained simulation yields
$r = 0.949$ vs.\ $0.931$ for Gaussian blur ($\Delta r = +0.018$;
Table~\ref{tab:ablation}, Fig.~\ref{fig:ablation}). A paired bootstrap
test on the same 75 subjects finds this difference is not statistically
significant (95\,\% CI on $\Delta r$: [$-0.021$, $+0.057$];
Wilcoxon $p = 0.847$). The result is a small positive trend rather
than a confirmed superiority; we retain the comparison as a
reproducible recipe baseline. These results are not directly comparable
to Table~\ref{tab:main} (different training conditions).

\begin{table}[!t]
\caption{Simulation ablation (separate training conditions; not comparable
to Table~\ref{tab:main})}
\label{tab:ablation}
\centering
\setlength{\tabcolsep}{5pt}
\begin{tabular}{lcc}
\toprule
Pre-training condition & Pearson $r$ [95\,\% CI] & MAE \\
\midrule
Physics-constrained sim (ours) & 0.949 [0.921--0.968] & 0.015 \\
Gaussian blur ($\sigma=2$)     & 0.931 [0.893--0.956] & 0.015 \\
\midrule
$\Delta$ (physics $-$ blur)    & $+0.018$ [$-0.021$, $+0.057$]$^*$ & $0.000$ \\
\bottomrule
\multicolumn{3}{l}{\footnotesize IXI pre-training + OASIS fine-tuning; $n_{\text{test}} = 75$
(IXI 80-train/75-test split, no domain-adaptive fine-tuning). Not comparable to Table~\ref{tab:main}.}\\
\multicolumn{3}{l}{\footnotesize $^*$Bootstrap 95\,\% CI; Wilcoxon $p = 0.847$ — difference not significant.}
\end{tabular}
\end{table}

\begin{figure}[!t]
\centering
\includegraphics[width=0.8\textwidth]{fig_ablation_gaussblur.png}
\caption{Physics-constrained simulation vs.\ Gaussian-blur pre-training
($n_{\text{test}} = 75$, separate training conditions from Table~\ref{tab:main}).}
\label{fig:ablation}
\end{figure}

\subsection{ViT3D vs.\ CNN3D on Real 64\,mT}
\label{sec:vit_cnn}

On real 64\,mT scans without adaptation ($n = 23$), CNN3D achieves higher
Pearson $r$ (0.514 vs.\ 0.291), while ViT3D achieves lower MAE (0.040
vs.\ 0.076). CNN3D's systematic bias of $-0.076$\,nWBV is $3.8\times$ the
reference threshold of 0.020 and renders it clinically unusable regardless of
its higher correlation. ViT3D's pre-adaptation mean bias of $+0.040$\,nWBV (which also equals
the unadapted MAE, since errors are near-uniformly signed) is corrected
to $-0.0022$ following LOOCV adaptation.

On OASIS-1, CNN3D outperforms ViT3D substantially (MAE $= 0.024$
vs.\ $0.058$; $r = 0.877$ vs.\ $0.722$). We attribute the reversal on real
64\,mT hardware to the narrow nWBV range (width 0.056 vs.\ 0.250 for
OASIS-1): CNN3D's larger output variance inflates $r$ while also inflating
MAE via systematic bias, whereas ViT3D's global self-attention produces more
stable absolute-value predictions in the narrow-range regime. A paired
bootstrap test on the $n=23$ subjects (100{,}000 resamples) confirms
ViT3D achieves lower absolute error in 20/23 subjects ($p < 0.001$).

\begin{figure}[!t]
\centering
\includegraphics[width=0.8\textwidth]{fig5_model_comparison.png}
\caption{Method comparison on real 64\,mT data ($n=23$): SynthSeg+,
CNN3D (no adaptation), and ViT3D (no adaptation).}
\label{fig:model_comparison}
\end{figure}

\begin{figure}[!t]
\centering
\includegraphics[width=0.8\textwidth]{fig6_baseline_comparison.png}
\caption{Four-way comparison on real 64\,mT data including LOOCV-adapted
ViT3D: Pearson $r$, MAE, and inference latency.}
\label{fig:baseline}
\end{figure}

\subsection{Cross-Session LOOCV Results}
\label{sec:loocv_results}

The primary evaluation uses 23-fold cross-session LOOCV. The adapter is
trained on session~HFC of 22 subjects and evaluated on session~HFE of the
held-out subject. Results:

\begin{itemize}
  \item MAE $= 0.0134$ [bootstrap 95\,\% CI: 0.010--0.017]; entire CI below
        the 0.020 reference threshold
  \item RMSE $= 0.016$, mean bias $= -0.0022$\,nWBV
  \item 19/23 subjects below the 0.020 threshold
  \item 4 subjects exceeding the threshold: HYPE21 (0.033), HYPE01 (0.029),
        HYPE16 (0.026), HYPE19 (0.021). Failures occur at both tails of
        the cohort range: HYPE16 (0.758) and HYPE21 (0.752) are at the
        low end (over-predicted); HYPE01 (0.806) and HYPE19 (0.801) are
        at the high end (under-predicted). Both patterns reflect regression
        toward the cohort mean
\end{itemize}

\textbf{Constant-mean baseline.}
The constant-mean LOOCV baseline (predict leave-one-out cohort mean
for each held-out subject) yields MAE $= 0.0126$. The ViT3D LOOCV MAE
($0.0134$) is therefore not meaningfully smaller than the
constant-mean MAE in absolute terms, which is an expected consequence
of the 0.056 nWBV range: in such a regime, the MAE metric is
dominated by the cohort spread rather than by the model. The result that distinguishes the model from a constant predictor is
the inter-session ICC reported below, which requires input-conditional
behaviour and cannot be produced by a constant baseline. Ground-truth
nWBV shows strong age-related decline ($\rho = -0.778$, $p < 0.001$),
confirming expected cohort biology; model predictions do not
significantly correlate with age ($\rho = +0.232$, $p = 0.287$),
indicating the model does not reliably capture this biological trend
at $n = 23$. The headline contribution of the LOOCV result is
therefore that the model's predictions remain bounded within the
healthy-cohort range on real 64\,mT hardware after adaptation, not
that they outperform a constant predictor on absolute error.

Fig.~\ref{fig:age} shows ground-truth and predicted nWBV plotted
against subject age side by side. Ground-truth nWBV declines strongly
with age ($\rho = -0.778$, $p < 0.001$), reflecting known cohort
biology. Model predictions do not significantly correlate with age
($\rho = +0.232$, $p = 0.287$), indicating that the model does not
reliably encode this trend at $n = 23$; this is an acknowledged
limitation of the current feasibility study.

\begin{figure}[!t]
\centering
\includegraphics[width=\textwidth]{fig4_age_nwbv.png}
\caption{(A) Ground-truth nWBV vs.\ age: strong decline ($\rho = -0.778$,
$p < 0.001$) reflects cohort biology. (B) LOOCV-adapted ViT3D predictions
vs.\ age: no significant correlation ($\rho = +0.232$, $p = 0.287$),
indicating the model does not reliably capture age-related decline at
$n = 23$. Both panels: real 64\,mT data, $n = 23$.}
\label{fig:age}
\end{figure}

Pearson $r$ is not reported as the primary metric for the 64\,mT LOOCV
evaluation. The ds006557 cohort nWBV range (0.752--0.808, width $= 0.056$)
is approximately 4.5$\times$ narrower than OASIS-1 (0.640--0.890, width
$\approx 0.250$). In such a restricted range, Pearson $r$ is structurally
sensitive to small prediction perturbations and is not indicative of
absolute prediction accuracy. MAE maps directly to the reference threshold
and is the appropriate primary measure~\cite{buckner2004}.

Fig.~\ref{fig:loocv_scatter} shows the predicted vs.\ true nWBV scatter
for all 23 held-out subjects under LOOCV.

\begin{figure}[!t]
\centering
\includegraphics[width=0.8\textwidth]{fig_loocv_scatter.png}
\caption{Cross-session LOOCV predicted vs.\ true nWBV ($n=23$).
Each point is one held-out subject. The shaded band shows $\pm$MAE.
The dashed line is the identity (perfect prediction).}
\label{fig:loocv_scatter}
\end{figure}

Bland--Altman analysis (Fig.~\ref{fig:ba}) confirms a near-zero mean bias
of $-0.0022$\,nWBV following adaptation.

\begin{figure}[!t]
\centering
\includegraphics[width=0.8\textwidth]{fig_bootstrap_ba.png}
\caption{(A) Bootstrap MAE distribution (10{,}000 resamples) for
LOOCV-adapted ViT3D on real 64\,mT ($n=23$); 95\,\% CI $= [0.010,\,0.017]$,
entirely below the reference threshold of 0.020\,nWBV.
(B) Bland--Altman plot: mean bias $= -0.0022$\,nWBV.}
\label{fig:ba}
\end{figure}

\subsection{Inter-Session Reliability: ICC(3,1)}

Using the same fold model to predict both HFE (primary) and HFC
(reliability check) of the held-out subject provides paired predictions
across two independent 64\,mT sessions. ICC(3,1) (two-way mixed-effects,
single measures, absolute agreement) $= 0.615$ [bootstrap 95\,\% CI:
0.236--0.866]: the point estimate is consistent with moderate
inter-session reliability by standard benchmarks~\cite{koo2016icc},
though the wide confidence interval [0.236--0.866] precludes a definitive
reliability classification at $n = 23$. The interval width is consistent
with the test-retest variability inherent in 64\,mT
hardware~\cite{vasa2025imaneur}.

\begin{figure}[!t]
\centering
\includegraphics[width=0.8\textwidth]{fig_loocv_icc.png}
\caption{Inter-session reliability. (Left) Per-subject absolute error
with the 0.020 reference threshold marked. (Right) Scatter of HFC vs.\
HFE LOOCV predictions ($n=23$); ICC(3,1) $= 0.615$
[95\,\% CI: 0.236--0.866].}
\label{fig:loocv_icc}
\end{figure}

\subsection{Simulated Dementia: CDR-Stratified Evaluation}
\label{sec:simulated_dementia}

No publicly available dataset provides real 64\,mT scans from dementia
patients. To characterize expected failure, we apply the physics-constrained
simulator to the OASIS-1 test set ($n = 38$, CDR 0--1.0) and evaluate
the OASIS fine-tuned ViT3D without domain adaptation (Table~\ref{tab:cdr_sim}).

\begin{table}[!t]
\caption{CDR-stratified MAE on physics-simulated 64\,mT inputs
(OASIS-1 test set, $n=38$; no domain adaptation applied)}
\label{tab:cdr_sim}
\centering
\setlength{\tabcolsep}{6pt}
\begin{tabular}{lccc}
\toprule
CDR group & $n$ & MAE $\pm$ SD & Interpretation \\
\midrule
No CDR label    & 20 & $0.016 \pm 0.010$ & Within threshold \\
CDR $= 0.0$     &  8 & $0.086 \pm 0.051$ & $4.4\times$ threshold \\
CDR $= 0.5$     &  8 & $0.099 \pm 0.031$ & $5.0\times$ threshold \\
CDR $= 1.0$     &  2 & $0.157 \pm 0.012$ & $7.9\times$ threshold \\
\midrule
Overall         & 38 & \multicolumn{2}{l}{MAE $= 0.055$, Pearson $r = 0.235$ ($p = 0.156$)} \\
\bottomrule
\multicolumn{4}{l}{\footnotesize Reference threshold: MAE $= 0.020$~\cite{buckner2004}.}\\
\multicolumn{4}{l}{\footnotesize CDR=1.0 MAE is $10\times$ the no-CDR-label group MAE (0.157 vs.\ 0.016).}\\
\multicolumn{4}{l}{\footnotesize CDR=1.0 $n=2$ makes this estimate directional rather than conclusive.}
\end{tabular}
\end{table}

The model predicts $\approx 0.837$ for most atrophied subjects regardless
of true nWBV, reflecting regression to the healthy-adult mean. The
$10\times$ MAE increase from no-CDR-label to CDR=1.0 (0.016 to 0.157)
indicates that atrophic brains are the dominant failure mode and that
the current model is \emph{not} appropriate for dementia patient
assessment, noting that CDR=1.0 $n=2$ makes this estimate directional
rather than conclusive.

\begin{figure}[!t]
\centering
\includegraphics[width=0.8\textwidth]{fig_simulated_dementia.png}
\caption{CDR-stratified MAE on physics-simulated 64\,mT inputs
(OASIS-1 test set, $n=38$; no domain adaptation). Error bars show $\pm$1\,SD.
The dashed line marks the 0.020 reference threshold. Only the no-CDR-label
group ($n=20$) falls within threshold.}
\label{fig:cdr_sim}
\end{figure}

\subsection{Pseudo-Label Ablation}
\label{sec:pseudo_ablation}

The LOOCV adaptation uses FastSurfer GT labels. In deployment, only
SynthSeg+ pseudo-labels would be available. We evaluate label-source
sensitivity on a fixed 15-train/8-test split:

\textbf{Label agreement (all 23 subjects):} SynthSeg+ vs.\ FastSurfer GT:
Pearson $r = 0.918$ ($p < 0.001$), MAE $= 0.005$, mean bias
$= -0.002$\,nWBV.

\begin{table}[!t]
\caption{Pseudo-label ablation: adapter test MAE by label source
(fixed 15-train/8-test split, seed 42)}
\label{tab:pseudo}
\centering
\setlength{\tabcolsep}{6pt}
\begin{tabular}{lccc}
\toprule
Label source & $n_{\text{train}}$ & MAE [95\,\% CI] & Bias \\
\midrule
FastSurfer GT (research)    & 15 & 0.012 [0.005--0.021] & $+0.006$ \\
SynthSeg+ pseudo (deployment) & 15 & 0.017 [0.012--0.023] & $-0.011$ \\
\midrule
MAE penalty & & $+0.005$ & \\
\bottomrule
\multicolumn{4}{l}{\footnotesize Test ground truth always FastSurfer GT. Both MAE point estimates below 0.020 threshold (CI upper bounds up to 0.023).}\\
\multicolumn{4}{l}{\footnotesize Pseudo-label penalty of $+0.005$ keeps both label sources within the feasibility threshold.}
\end{tabular}
\end{table}

Both label sources yield test MAE below the 0.020 reference threshold.
The $+0.005$ penalty from pseudo-labels indicates that the adapter
training stage does not require a FastSurfer installation, which is a
practical requirement for deployment in environments where high-field
ground truth is unavailable.

\begin{figure}[!t]
\centering
\includegraphics[width=0.8\textwidth]{fig_pseudolabel_ablation.png}
\caption{Pseudo-label ablation. (Left) SynthSeg+ vs.\ FastSurfer GT label
agreement across all 23 subjects (Pearson $r = 0.918$, $p < 0.001$).
(Right) Adapter test MAE for both label sources; both remain below the
0.020 reference threshold.}
\label{fig:pseudolabel}
\end{figure}

\subsection{Uncertainty Quantification}
\label{sec:uncertainty}

Predictive accuracy and uncertainty calibration are conceptually
distinct properties of a regression model and are reported separately
here. The accuracy result --- LOOCV MAE $= 0.0134$\,nWBV with the
entire bootstrap interval below the 0.020 reference threshold --- does
not depend on the calibration of any associated uncertainty estimator.

We additionally evaluated the most commonly applied lightweight
uncertainty estimator, MC Dropout~\cite{gal2016dropout}, with $N = 100$
stochastic forward passes per scan. The mean 95\,\% interval width is
$0.029$\,nWBV on both the OASIS-1 test set (Fig.~\ref{fig:ci}) and the
real 64\,mT cohort (Fig.~\ref{fig:ci_real}). Empirical coverage at the
nominal 95\,\% level is 4.3\,\% (Fig.~\ref{fig:calibration}), and
interval width is not significantly correlated with absolute prediction error
($r = +0.34$, $p = 0.108$, $n = 23$). The primary evidence of miscalibration
is the 4.3\,\% empirical coverage against a nominal 95\,\% target; the positive
but non-significant $r$ value is consistent with the direction one would expect
from a well-calibrated model and does not by itself indicate miscalibration.
Both observations are consistent with the
known failure mode of dropout-based uncertainty under
covariate shift~\cite{ovadia2019cancalibration}: the variational
approximation underlying MC Dropout assumes that the test distribution
matches the training distribution, which does not hold when the model
is trained partly on physics-simulated 64\,mT data and tested on real
hardware.

This calibration result is reported for completeness and is not used
to qualify any of the accuracy claims in this paper. Calibrated
predictive intervals require either (i) a held-out real-hardware
calibration set with sufficient pathology coverage, or (ii) split
conformal prediction~\cite{angelopoulos2021conformal} computed on
prospective acquisitions. Neither is feasible within the present
ds006557 cohort: setting aside calibration data from $n = 23$ would
shrink the LOOCV evaluation set below the threshold for meaningful
bootstrap inference. Conformal calibration is therefore deferred to
the prospective multi-site validation already identified as the next
step (Section~\ref{sec:limitations}). MC Dropout intervals must
\emph{not} be used for individual-level clinical decisions on the
basis of this paper.

\begin{figure}[!t]
\centering
\includegraphics[width=\textwidth]{fig8a_ci_oasis.png}
\caption{Per-subject nWBV predictions with MC Dropout 95\,\% CI ---
OASIS-1 test set ($n=38$). Empirical coverage $= 4.3\%$ (nominal $= 95\%$);
CIs are severely underestimated and not suitable for clinical use.}
\label{fig:ci}
\end{figure}

\begin{figure}[!t]
\centering
\includegraphics[width=\textwidth]{fig8b_ci_real64mt.png}
\caption{MC Dropout uncertainty --- real 64\,mT data ($n=23$, $N=100$
passes). CI width is not significantly correlated with error
($r = +0.34$, $p = 0.108$); the 4.3\,\% empirical coverage is the
primary indicator of miscalibration under domain shift.}
\label{fig:ci_real}
\end{figure}

\subsection{Failure Analysis}
\label{sec:failure}

\begin{figure}[!t]
\centering
\includegraphics[width=\textwidth]{fig9_failure_analysis.png}
\caption{Failure analysis: (A) residuals by true nWBV, (B) MAE by nWBV
tertile, (C) error vs.\ age on real 64\,mT LOOCV, (D) MC Dropout CI
width vs.\ absolute error.}
\label{fig:failure}
\end{figure}

\textbf{(A) Residuals.} Large residuals ($> 0.050$) occur exclusively at
nWBV $< 0.76$. The model predicts $\approx 0.837$ for severely atrophied
brains, reflecting regression to the healthy-adult mean due to
under-representation of atrophic subjects in training.

\textbf{(B) MAE by nWBV tertile.} Lowest tertile (atrophy range,
nWBV $< 0.76$): MAE $= 0.125$. Highest tertile (normal range): MAE $= 0.009$.
Ratio: $14\times$. All subjects exceeding the 0.020 reference threshold are
in the lowest tertile.

\textbf{(C) LOOCV failures.} The 4 subjects exceeding 0.020 on real 64\,mT
LOOCV (HYPE21: 0.033, HYPE01: 0.029, HYPE16: 0.026, HYPE19: 0.021) show
failures at \emph{both} tails of the cohort nWBV range: HYPE21 (0.752)
and HYPE16 (0.758) are at the low end and are over-predicted; HYPE01
(0.806) and HYPE19 (0.801) are at the high end and are under-predicted.
Both patterns are consistent with regression toward the cohort mean ---
the model compresses predictions toward the centre of its training range.

\textbf{(D) Calibration.} MC Dropout CI width is not significantly correlated
with absolute error ($r = +0.34$, $p = 0.108$, $n = 23$). The decisive
evidence of miscalibration is the reliability curve
(Fig.~\ref{fig:calibration}): nominal 95\,\% achieves only 4.3\,\% empirical
coverage, a severe systematic underestimation that confirms MC Dropout
intervals are not suitable for deployment use under this domain shift.

Table~\ref{tab:worst5} lists the five highest-error subjects from the
LOOCV evaluation. Failures cluster at the edges of the cohort range:
subjects with low nWBV are over-predicted, while subjects with high nWBV
are under-predicted --- both driven by regression toward the cohort mean.

\begin{table}[!t]
\caption{Worst-5 LOOCV subjects by absolute error (real 64\,mT, $n=23$).
Failures occur at both tails of the cohort range (low nWBV: over-predicted;
high nWBV: under-predicted), consistent with regression toward the cohort mean.}
\label{tab:worst5}
\centering
\setlength{\tabcolsep}{8pt}
\renewcommand{\arraystretch}{1.1}
\begin{tabular}{lcccc}
\toprule
Subject & True nWBV & Predicted nWBV & $|$Error$|$ & Threshold exceeded? \\
\midrule
sub-HYPE21 & 0.752 & 0.786 & 0.033 & \textbf{Yes} ($1.7\times$) \\
sub-HYPE01 & 0.806 & 0.777 & 0.029 & \textbf{Yes} ($1.5\times$) \\
sub-HYPE16 & 0.758 & 0.784 & 0.026 & \textbf{Yes} ($1.3\times$) \\
sub-HYPE19 & 0.801 & 0.780 & 0.021 & \textbf{Yes} ($1.0\times$) \\
sub-HYPE02 & 0.794 & 0.775 & 0.019 & No ($0.95\times$) \\
\bottomrule
\multicolumn{5}{l}{\footnotesize Threshold: 0.020\,nWBV~\cite{buckner2004}. $N\times$ = ratio to threshold.}\\
\multicolumn{5}{l}{\footnotesize Low-nWBV subjects (HYPE21, HYPE16) are over-predicted; high-nWBV subjects}\\
\multicolumn{5}{l}{\footnotesize (HYPE01, HYPE19) are under-predicted. Both patterns reflect regression toward cohort mean.}
\end{tabular}
\end{table}

\begin{figure}[!t]
\centering
\includegraphics[width=\textwidth]{fig_calibration.png}
\caption{MC Dropout calibration analysis ($n=23$, real 64\,mT).
(A) Reliability curve: empirical coverage vs.\ nominal confidence level.
The model is severely overconfident --- intervals are far too narrow,
with nominal 95\,\% achieving only 4.3\,\% empirical coverage.
(B) Calibration error at key confidence levels. MC Dropout intervals
are reported for transparency only and must not be used for clinical
decisions.}
\label{fig:calibration}
\end{figure}

% ============================================================
\section{Discussion}
\label{sec:discussion}
% ============================================================

\subsection{Clinical Feasibility and Scope}

Following lightweight session-level adaptation (769 parameters,
$< 1$\,min training), ViT3D produces nWBV estimates on real 64\,mT
hardware with MAE $= 0.0134$\,nWBV [0.010--0.017] --- approximately
19\,mL of brain parenchyma --- with the entire bootstrap interval below
the 0.020\,nWBV clinical reference threshold~\cite{buckner2004}.
Inference completes in 47\,ms on a standard GPU without FreeSurfer,
SynthSeg+, or any segmentation step (edge-device latency not yet
benchmarked). Ground-truth nWBV declines with age ($\rho = -0.778$, $p < 0.001$),
confirming expected cohort biology; model predictions do not
significantly correlate with age ($\rho = +0.232$, $p = 0.287$),
an acknowledged limitation at this sample size.

These properties support a feasibility claim in the specific setting of
compute-constrained monitoring of healthy adults (e.g.\ bedside or
low-resource environments), where software-heavy pipelines such as
FreeSurfer or SynthSeg+ are unavailable. Latency is benchmarked at
47\,ms on a standard GPU; edge-device latency is not yet measured and
is identified as future work. The claim is bounded by
three factors that we treat in detail below: validation is confined to a
single healthy cohort ($n = 23$), inter-session reliability has only a
preliminary point estimate with a wide bootstrap interval at this sample
size, and the model degrades sharply on atrophic brains. Clinical deployment
in dementia screening or diagnostic decision support is outside the
present scope.

\subsection{Architecture Choice: ViT3D vs.\ CNN3D}
\label{sec:vit_cnn_disc}

The architecture comparison is regime-specific, not absolute: the
superiority of ViT3D in the 64\,mT setting does not generalise to
3\,T or other wider-range regimes. The case for the choice rests on
three points of justification.

\textbf{Why MAE is the primary metric.}
The clinical reference threshold for distinguishing normal from
pathological nWBV change is an absolute quantity ($0.020$\,nWBV;
\cite{buckner2004}), defined in nWBV units rather than as a correlation
coefficient. MAE measures the same quantity that the threshold
constrains, so it is the metric that determines whether a prediction
falls inside or outside the acceptable error budget. Pearson $r$
measures the linear monotonic association between predictions and
targets and is invariant to systematic bias; in the narrow ds006557
range (nWBV width $0.056$), small absolute errors that exceed the
reference threshold can still yield arbitrarily high $r$. We therefore
pre-specified MAE as the primary metric before any model evaluation.

\textbf{Why ViT3D's adapter correction matters.}
On real 64\,mT hardware without adaptation, unadapted ViT3D achieves
lower MAE than unadapted CNN3D (0.040 vs.\ 0.076), establishing a
like-for-like advantage before any adaptation is applied. After a
769-parameter session-level adapter (final LayerNorm and regression
head), ViT3D's mean bias is further reduced from $+0.040$ to $-0.0022$\,nWBV
and LOOCV MAE collapses to $0.0134$ [0.010--0.017]. Whether the same
adapter would correct CNN3D's $-0.076$ bias to a comparable degree
was not tested and is identified as future work; the architecture
comparison here is therefore based on the unadapted like-for-like
result (ViT3D MAE 0.040 vs.\ CNN3D MAE 0.076) and on the global
self-attention argument below, not on a claim that CNN3D's bias is
structurally uncorrigible.

\textbf{Domain-specific, not general.}
We do not claim Transformer-class architectures dominate convolutional
ones in medical imaging. CNN3D is materially better on OASIS-1
(wide-range, 3\,T-derived inputs), and that result is reported in
Table~\ref{tab:main} without qualification. The advantage observed
here is specific to the narrow-range, high-noise, low-data 64\,mT
regime, where the corrigibility of the residual bias becomes the
dominant criterion. This finding is consistent with prior reports
that Vision Transformers exhibit lower texture bias and stronger
global-shape sensitivity than CNNs~\cite{dosovitskiy2020image}; both
are advantageous when the target is a global volumetric ratio and the
local texture is degraded by Rician noise.

\subsection{SynthSeg+ as the Reference Standard}

SynthSeg+ ($r = 0.918$, MAE $= 0.005$) sets the achievable accuracy
ceiling for nWBV estimation on this cohort. ViT3D does not match this
accuracy, but the relevant comparison is deployment context: SynthSeg+
requires a 12\,GB FreeSurfer installation and 2--3\,min per scan,
whereas ViT3D requires only PyTorch and 47\,ms of inference on a
standard GPU. In compute-constrained settings where SynthSeg+ cannot
be co-located with the scanner, ViT3D's accuracy--latency trade-off
becomes the operative
constraint rather than a limitation.

\subsection{Inter-Session Reliability}
\label{sec:discussion_icc}

The ICC(3,1) point estimate of $0.615$ is consistent with moderate
measurement consistency by the benchmarks of Koo and
Li~\cite{koo2016icc}, but the bootstrap confidence interval
[0.236--0.866] precludes a definitive reliability classification at
$n = 23$. Two factors contribute jointly to the point estimate. The
adapter is applied independently to the HFC and HFE sessions, so any
session-to-session instability in the pre-adaptation features
propagates into the predictions. In addition, V{\'a}{\v s}a et
al.~\cite{vasa2025imaneur} characterised 64\,mT hardware as having
non-negligible native test-retest variability, which sets an upper
bound on the ICC any model can achieve on this hardware. The reported
value therefore reflects the joint reliability of the model and the
scanner. The current interval width means the reliability claim should
be interpreted as directional; the lower bound of 0.236 falls in the
poor-to-fair range and cannot be ruled out at this sample size. The
wide bootstrap interval is a sample-size effect at $n = 23$ rather
than evidence of instability and is expected to narrow with larger
prospective cohorts. At this level of evidence the model is
appropriate for group-level or longitudinal research use; individual
clinical decisions would require either a larger reliability study or
a calibrated rejection mechanism not pursued here.

\subsection{Limitations}
\label{sec:limitations}

We separate methodological limitations of the present study from
clinical scope conditions that govern future use.

\subsubsection*{Methodological limitations}

\textbf{Sample size and dataset availability.}
Real-hardware evaluation uses ds006557 ($n = 23$), the only public
paired 64\,mT/3\,T brain MRI dataset at the time of writing. The
sample size therefore reflects the state of publicly available ULF
data rather than a study-design choice. Bootstrap intervals are
reported throughout and the cross-session LOOCV protocol is fully
subject-independent.

\textbf{Physics-simulation acquisition gap.}
The simulator models tissue-specific T1/T2 relaxation, Rician noise,
and B0 inhomogeneity, but does not model the full Hyperfine
acquisition chain (ETL-80 averaging, $k$-space sampling pattern,
compressed-sensing reconstruction). The resulting SNR mismatch
produces a pre-adaptation bias of $+0.040$\,nWBV, which the
769-parameter adapter reduces to $-0.0022$. Relative tissue contrast
ratios appear to be preserved and to drive the regression; modelling
the full acquisition chain is a tractable refinement.

\textbf{Uncertainty calibration.}
MC Dropout intervals are severely under-covered under domain shift
(Section~\ref{sec:uncertainty}), reflecting a known failure mode of
dropout-based variational uncertainty rather than a property of the
underlying point predictions. These intervals are reported for
completeness and are not used to qualify any accuracy result. Split
conformal calibration on prospective acquisitions is the appropriate
remedy and is identified as future work.

\subsubsection*{Clinical scope conditions}

\textbf{Healthy-only validation cohort.}
ds006557 contains only healthy adults. The CDR-stratified
physics-simulated evaluation (Section~\ref{sec:simulated_dementia})
quantifies the expected degradation on pathological brains: CDR=1.0
subjects exhibit MAE $0.157$, an order of magnitude above the
healthy-reference MAE ($0.016$), noting that CDR=1.0 $n=2$ makes
this estimate directional rather than conclusive. The model is
therefore not appropriate for use on atrophic or pathological cohorts
in its current form. Prospective acquisition of paired 64\,mT/3\,T
data from CDR-rated patients is the necessary next step.

\textbf{Regression-to-the-mean failure on low-nWBV inputs.}
MAE in the lowest OASIS-1 nWBV tertile ($< 0.76$) is $0.125$, a
$14\times$ increase over the normal-range tertile. This is driven by
under-representation of atrophic subjects in OASIS-1 training rather
than by the architecture or the simulator. Retraining on
atrophy-enriched cohorts such as OASIS-3 ($n > 1{,}000$, longitudinal
CDR) is the path to addressing this failure mode.

\textbf{Inter-session reliability.}
The ICC(3,1) reliability analysis and its clinical scope implications
are discussed in Section~\ref{sec:discussion_icc}.

\textbf{Single-site evaluation.}
All real-hardware data come from one Hyperfine Swoop generation at a
single site. Multi-site, multi-generation reproducibility is not
established and is identified as a prerequisite for any deployment.

\subsection{Future Work}

(i) Retraining on OASIS-3 to improve atrophy-range coverage;
(ii) Modelling the full ETL acquisition chain in simulation to close the SNR gap;
(iii) Prospective paired 64\,mT/3\,T acquisition from CDR-rated patients;
(iv) Three-way adapter strategy ablation (head-only vs.\ LN+head vs.\ full fine-tune);
(v) Conformal prediction for calibrated uncertainty quantification;
(vi) Multi-site validation across Hyperfine scanner generations.

% ============================================================
\section{Conclusion}
\label{sec:conclusion}
% ============================================================

A compact 3D Vision Transformer (ViT3D, 4.23\,M parameters) trained on
physics-constrained 64\,mT simulations and fine-tuned on OASIS-1
estimates normalized whole-brain volume from bedside Hyperfine Swoop
MRI in 47\,ms (standard GPU) without segmentation software. Under fully
subject-independent cross-session LOOCV ($n = 23$), MAE
$= 0.0134$\,nWBV [0.010--0.017] --- entirely below the 0.020 clinical
reference threshold, but comparable to a constant-mean LOOCV baseline
(MAE $= 0.0126$) given the narrow ds006557 nWBV range. Discriminative
evidence beyond MAE is provided by ICC(3,1) $= 0.615$ [0.236--0.866], a preliminary inter-session
reliability estimate whose wide confidence interval precludes a
definitive classification at this sample size.
Physics-constrained simulation
shows a small positive trend over Gaussian-blur degradation
($\Delta r = +0.018$, not statistically significant at $n=75$,
bootstrap 95\,\% CI [$-0.021$, $+0.057$]), and SynthSeg+ pseudo-labels are a viable
deployment substitute for FastSurfer ground truth (MAE penalty
$+0.005$).

The results are bounded to healthy adults. In physics-simulated
dementia conditions, CDR=1.0 subjects exhibit MAE an order of magnitude
above the healthy reference, identifying atrophic brains as the
principal unresolved failure mode. The contribution of this work is a
reproducible feasibility baseline for automated brain morphometry from
ultra-low-field bedside MRI, together with a precise characterisation
of the conditions --- atrophy-enriched retraining and prospective
clinical validation --- under which the approach can be extended toward
diagnostic use.

% ============================================================
\appendix
\section{Training Hyperparameters}
\label{app:hyperparams}

\begin{table}[!h]
\caption{Hyperparameters for reproducibility}
\label{tab:hyper}
\centering
\setlength{\tabcolsep}{6pt}
\begin{tabular}{lll}
\toprule
Parameter & Stage 2 (OASIS fine-tune) & Stage 3 (LOOCV adapt.) \\
\midrule
Optimizer             & AdamW              & AdamW \\
Learning rate         & $5\times10^{-5}$   & $1\times10^{-4}$ \\
Weight decay          & $10^{-4}$          & $10^{-4}$ \\
Max epochs            & 50                 & 80 \\
Early stop patience   & 10 (val Pearson $r$) & --- \\
Batch size            & 4                  & 4 \\
LR schedule           & Cosine annealing   & Cosine annealing \\
Augmentations/scan    & 2                  & 5 \\
Trainable parameters  & 4.23M (all)        & 769 (LayerNorm + head) \\
Input size            & $64^3$ voxels      & $64^3$ voxels \\
Intensity normalisation & 1--99 percentile clip & 1--99 percentile clip \\
Random seed           & 42                 & 42 \\
LOOCV folds           & ---                & 23 \\
Bootstrap resamples   & ---                & 10{,}000 \\
\bottomrule
\end{tabular}
\end{table}

% ============================================================
\section*{Acknowledgment}

The authors acknowledge the OASIS project (P50 AG05681, P01 AG03991,
R01 AG021910) and OpenNeuro for data access.

% ============================================================
\begin{thebibliography}{00}

\bibitem{who2023}
World Health Organization, ``Dementia,'' WHO Fact Sheet, Sep. 2023.

\bibitem{fotenos2005}
A.~F. Fotenos, A.~Z. Snyder, L.~E. Girton, J.~C. Morris, and R.~L. Buckner,
``Normative estimates of cross-sectional and longitudinal brain volume
decline in aging and AD,''
\emph{Neurology}, vol.~64, no.~6, pp.~1032--1039, 2005.

\bibitem{buckner2004}
R.~L. Buckner \emph{et al.},
``Molecular, structural, and functional characterization of Alzheimer's
disease,''
\emph{J. Neurosci.}, vol.~25, no.~34, pp.~7709--7717, 2005.

\bibitem{jack2018nia}
C.~R. Jack \emph{et al.},
``NIA-AA Research Framework: Toward a biological definition of
Alzheimer's disease,''
\emph{Alzheimers Dement.}, vol.~14, no.~4, pp.~535--562, 2018.

\bibitem{fischl2012freesurfer}
B.~Fischl, ``FreeSurfer,''
\emph{NeuroImage}, vol.~62, no.~2, pp.~774--781, 2012.

\bibitem{smith2002fsl}
S.~M. Smith \emph{et al.},
``Advances in functional and structural MR image analysis and
implementation as FSL,''
\emph{NeuroImage}, vol.~23, suppl.~1, pp.~S208--S219, 2004.

\bibitem{prince2016world}
M.~Prince \emph{et al.},
``World Alzheimer Report 2016,''
Alzheimer's Disease International, London, 2016.

\bibitem{hyperfine2021}
Hyperfine Inc., ``Swoop Portable MR Imaging System: 510(k) Summary,''
FDA K192972, 2020.

\bibitem{kimberly2023bedside}
W.~T. Kimberly \emph{et al.},
``Portable low-field magnetic resonance imaging enables brain imaging
at the bedside,''
\emph{npj Digit. Med.}, vol.~6, no.~1, 2023.

\bibitem{vasa2025}
F.~V{\'a}{\v s}a \emph{et al.},
``Longitudinal paired 64\,mT brain MRI dataset (ds006557),''
OpenNeuro, 2025.

\bibitem{iglesias2023synthsr}
J.~E. Iglesias \emph{et al.},
``SynthSR: A public AI tool to turn heterogeneous clinical brain scans
into high-resolution T1-weighted images for 3D morphometry,''
\emph{Sci. Adv.}, vol.~9, no.~5, 2023, Art.~no.~eadd3607.

\bibitem{zhao2021superresolution}
C.~Zhao, A.~Carass, B.~E. Dewey, and J.~L. Prince,
``Self super-resolution for magnetic resonance images using deep
networks,''
in \emph{Proc. IEEE ISBI}, 2020, pp.~365--368.

\bibitem{chen2019synergistic}
C.~Chen, Q.~Dou, H.~Chen, J.~Qin, and P.~A. Heng,
``Synergistic image and feature adaptation: Towards cross-modality
domain adaptation for medical image segmentation,''
in \emph{Proc. AAAI}, vol.~33, 2019, pp.~865--872.

\bibitem{billot2023synthseg}
B.~Billot \emph{et al.},
``SynthSeg: Segmentation of brain MRI scans of any contrast and
resolution without retraining,''
\emph{Med. Image Anal.}, vol.~86, 2023, Art.~no.~102789.

\bibitem{dosovitskiy2020image}
A.~Dosovitskiy \emph{et al.},
``An image is worth 16$\times$16 words: Transformers for image
recognition at scale,''
in \emph{Proc. ICLR}, 2021.

\bibitem{chen2021transunet}
J.~Chen \emph{et al.},
``TransUNet: Transformers make strong encoders for medical image
segmentation,''
\emph{arXiv:2102.04306}, 2021.

\bibitem{hatamizadeh2021swin}
A.~Hatamizadeh \emph{et al.},
``Swin UNETR: Swin transformers for semantic segmentation of brain
tumors in MRI images,''
in \emph{Proc. MICCAI BraTS}, 2021, pp.~272--284.

\bibitem{cole2017predicting}
J.~H. Cole \emph{et al.},
``Predicting brain age with deep learning from raw imaging data,''
\emph{NeuroImage}, vol.~163, pp.~115--124, 2017.

\bibitem{jonsson2019brain}
B.~A. Jonsson \emph{et al.},
``Brain age prediction using deep learning uncovers associated
sequence variants,''
\emph{Nat. Commun.}, vol.~10, Art.~no.~5409, 2019.

\bibitem{sheth2021assessmentultralow}
K.~N. Sheth \emph{et al.},
``Assessment of brain injury using portable, low-field magnetic resonance
imaging at the bedside,''
\emph{JAMA Neurol.}, vol.~78, no.~1, pp.~41--47, 2021.

\bibitem{ixi2017}
IXI Dataset, Brain Development, Imperial College London, 2017.
[Online]. Available: \url{https://brain-development.org/ixi-dataset/}

\bibitem{marcus2007oasis}
D.~S. Marcus \emph{et al.},
``Open Access Series of Imaging Studies (OASIS): Cross-sectional MRI
data in young, middle aged, nondemented, and demented older adults,''
\emph{J. Cogn. Neurosci.}, vol.~19, no.~9, pp.~1498--1507, 2007.

\bibitem{henschel2020fastsurfer}
L.~Henschel \emph{et al.},
``FastSurfer -- A fast and accurate deep learning based neuroimaging
pipeline,''
\emph{NeuroImage}, vol.~219, 2020, Art.~no.~117012.

\bibitem{deichmann2002}
R.~Deichmann, C.~Schwarzbauer, and R.~Turner,
``Optimisation of the 3D MDEFT sequence for anatomical brain imaging,''
\emph{NeuroImage}, vol.~21, no.~2, pp.~757--767, 2004.

\bibitem{peng2021brain}
H.~Peng, W.~Gong, C.~F. Beckmann, A.~Vedaldi, and S.~M. Smith,
``Accurate brain age prediction with lightweight deep neural networks,''
\emph{Med. Image Anal.}, vol.~68, Art.~no.~101871, 2021.

\bibitem{roy2019quicknat}
A.~G. Roy, S.~Conjeti, N.~Navab, and C.~Wachinger,
``QuickNAT: A fully convolutional network for quick and accurate
segmentation of neuroanatomy,''
\emph{NeuroImage}, vol.~186, pp.~713--727, 2019.

\bibitem{karani2018lifelong}
N.~Karani, K.~Chaitanya, C.~Baumgartner, and E.~Konukoglu,
``A lifelong learning approach to brain MR segmentation across scanners
and protocols,''
in \emph{Proc. MICCAI}, Granada, Spain, 2018, pp.~476--484.

\bibitem{dewey2019deepharmony}
B.~E. Dewey \emph{et al.},
``DeepHarmony: A deep learning approach to contrast harmonization across
scanner changes,''
\emph{Magn. Reson. Imaging}, vol.~64, pp.~160--170, 2019.

\bibitem{vasa2025imaneur}
F.~V{\'a}{\v s}a \emph{et al.},
``Ultra-low-field brain MRI morphometry: Test-retest reliability and
correspondence to high-field MRI,''
\emph{Imaging Neuroscience}, Art.~no.~a930, 2024.

\bibitem{arnold2022simulated}
T.~C. Arnold, S.~N. Baldassano, B.~Litt, and J.~M. Stein,
``Simulated diagnostic performance of low-field MRI: Harnessing
open-access datasets to evaluate novel devices,''
\emph{Magn. Reson. Imaging}, vol.~87, pp.~67--76, Apr. 2022.

\bibitem{rebsamen2020cortical}
M.~Rebsamen, C.~Rummel, M.~Reyes, R.~Wiest, and R.~McKinley,
``Direct cortical thickness estimation using deep learning-based anatomy
segmentation and cortex parcellation,''
\emph{Hum. Brain Mapp.}, vol.~41, no.~17, pp.~4804--4814, 2020.

\bibitem{islam2025aibrain}
K.~T. Islam \emph{et al.},
``AI improves consistency in regional brain volumes measured in
ultra-low-field MRI and 3T MRI,''
\emph{Front. Neuroimaging}, vol.~4, Art.~no.~1588487, Jun. 2025.

\bibitem{javadi2025diffusion}
M.~Javadi \emph{et al.},
``In-silico comparison of a diffusion model with conventionally trained
deep networks for translating 64\,mT to 3T brain FLAIR,''
\emph{Sci. Rep.}, vol.~15, Art.~no.~21806, 2025.

\bibitem{gopinath2025cortical}
K.~Gopinath \emph{et al.},
``From low field to high value: Robust cortical mapping from
low-field MRI,''
arXiv preprint arXiv:2505.12228, May 2025.

\bibitem{minocha2024lfmri}
A.~Minocha and S.~Yang,
``Low field MRI deep learning framework for non-hospitalizing early
detection and characterization of Alzheimer's disease pathology,''
in \emph{Proc. IEEE Int. Conf. Big Data}, 2024, pp.~1--8.

\bibitem{dayarathna2024adversarial}
S.~Dayarathna, K.~T. Islam, and Z.~Chen,
``Ultra low-field to high-field MRI translation using adversarial
diffusion,''
in \emph{Proc. IEEE ISBI}, 2024, pp.~1--4.

\bibitem{baljer2024paediatric}
L.~Baljer \emph{et al.},
``Ultra-low-field paediatric MRI in low- and middle-income countries:
Super-resolution using a multi-orientation {U-Net},''
\emph{Hum. Brain Mapp.}, vol.~46, Art.~no.~e70112, 2025.

\bibitem{koo2016icc}
T.~K. Koo and M.~Y. Li,
``A guideline of selecting and reporting intraclass correlation
coefficients for reliability research,''
\emph{J. Chiropr. Med.}, vol.~15, no.~2, pp.~155--163, 2016.

\bibitem{gal2016dropout}
Y.~Gal and Z.~Ghahramani,
``Dropout as a Bayesian approximation: Representing model uncertainty
in deep learning,''
in \emph{Proc. ICML}, 2016, pp.~1050--1059.

\bibitem{ovadia2019cancalibration}
Y.~Ovadia \emph{et al.},
``Can you trust your model's uncertainty? Evaluating predictive
uncertainty under dataset shift,''
in \emph{Proc. NeurIPS}, 2019, pp.~13991--14002.

\bibitem{angelopoulos2021conformal}
A.~N. Angelopoulos and S.~Bates,
``A gentle introduction to conformal prediction and distribution-free
uncertainty quantification,''
\emph{arXiv:2107.07511}, 2021.

\end{thebibliography}

\end{document}
